\chapter{Descrizione delle Strutture Dati}

\section{KeyValue}

La struttura \texttt{KeyValue} rappresenta la coppia chiave-valore memorizzata 
nella lista di collisione di ciascun nodo. È definita come segue:

\begin{lstlisting}
struct KeyValue {
    string key;
    string value;
};
\end{lstlisting}

Entrambi i campi sono di tipo \texttt{std::string}. La separazione della coppia 
in una struttura dedicata rende il codice più leggibile e facilita la gestione 
delle collisioni nella lista.

\section{RBNode}

\texttt{RBNode} rappresenta il nodo dell'albero Red-Black. Ogni nodo contiene:

\begin{itemize}
    \item \texttt{int id} --- identificatore numerico del nodo, ottenuto applicando 
    la funzione di hash alla chiave inserita;
    \item \texttt{Color color} --- colore del nodo, \texttt{RED} o \texttt{BLACK}, 
    necessario per mantenere le proprietà dell'albero Red-Black;
    \item \texttt{std::list<KeyValue> items} --- lista concatenata delle coppie 
    \texttt{(chiave, valore)} che condividono lo stesso valore di hash (gestione 
    delle collisioni);
    \item \texttt{RBNode* left}, \texttt{RBNode* right}, \texttt{RBNode* parent} 
    --- puntatori al figlio sinistro, figlio destro e nodo padre.
\end{itemize}

Il tipo enumerato \texttt{Color} è definito come:

\begin{lstlisting}
enum Color { RED, BLACK };
\end{lstlisting}

Il costruttore di \texttt{RBNode} inizializza il colore a \texttt{RED} 
(ogni nuovo nodo inserito è sempre rosso), aggiunge la coppia 
\texttt{(chiave, valore)} alla lista \texttt{items} e imposta tutti i 
puntatori a \texttt{nullptr}:

\begin{lstlisting}
RBNode(int id, const string& key, const string& value)
    : id(id), color(RED), left(nullptr), right(nullptr), parent(nullptr) {
    items.push_back({key, value});
}
\end{lstlisting}

\section{HashRBTree}

\texttt{HashRBTree} è la classe principale che incapsula l'intera struttura dati. 
Mantiene internamente due puntatori:

\begin{itemize}
    \item \texttt{RBNode* root} --- puntatore alla radice dell'albero;
    \item \texttt{RBNode* NIL} --- nodo sentinella condiviso da tutti i nodi 
    foglia e dalla radice prima di qualsiasi inserimento.
\end{itemize}

\subsection{Nodo Sentinella NIL}

Invece di utilizzare \texttt{nullptr} per rappresentare i nodi foglia e i 
puntatori mancanti, si adotta un nodo sentinella \texttt{NIL} con le seguenti 
caratteristiche:

\begin{itemize}
    \item colore \texttt{BLACK};
    \item tutti i puntatori (\texttt{left}, \texttt{right}, \texttt{parent}) 
    puntano a se stesso.
\end{itemize}

Questo approccio semplifica notevolmente l'implementazione di \texttt{insertFix} 
e \texttt{deleteFix}, eliminando la necessità di controllare \texttt{nullptr} 
ad ogni accesso ai campi del nodo: è sempre possibile accedere a 
\texttt{->color}, \texttt{->parent} ecc. senza rischio di dereferenziazione nulla.

\subsection{Interfaccia Pubblica}

La classe espone quattro metodi pubblici:

\begin{itemize}
    \item \texttt{void insert(const string\& key, const string\& value)} --- 
    inserisce una coppia \texttt{(chiave, valore)};
    \item \texttt{void search(const string\& key, const string\& value)} --- 
    ricerca una chiave e stampa il valore associato;
    \item \texttt{void remove(const string\& key, const string\& value)} --- 
    rimuove una coppia \texttt{(chiave, valore)};
    \item \texttt{void print()} --- stampa l'albero in forma grafica.
\end{itemize}

\subsection{Metodi Privati}

I metodi privati supportano le operazioni pubbliche:

\begin{itemize}
    \item \texttt{hashFunction} --- calcola l'hash della chiave;
    \item \texttt{leftRotate}, \texttt{rightRotate} --- rotazioni per il 
    bilanciamento;
    \item \texttt{insertFix}, \texttt{deleteFix} --- ripristino delle proprietà 
    RBT dopo inserimento e cancellazione;
    \item \texttt{transplant} --- sostituzione di un sottoalbero con un altro;
    \item \texttt{minimum} --- nodo con id minimo in un sottoalbero;
    \item \texttt{searchNode} --- ricerca interna per id;
    \item \texttt{printHelper}, \texttt{destroyTree} --- supporto a stampa e 
    distruzione.
\end{itemize}